\section{তথ্য ও যোগাযোগ প্রযুক্তি ব্যবহারে নৈতিকতা}
\label{sec:ch1-ethics}
\subsection{ধারণা ও মেধাসম্পদ}
\definition{ICT নৈতিকতা}{প্রযুক্তি ব্যবহারে ব্যক্তি, প্রতিষ্ঠান ও সমাজের অধিকার, নিরাপত্তা, গোপনীয়তা, সত্যতা ও দায়িত্ব সম্মান করার নীতিমালাকে ICT নৈতিকতা বলে।}
দ্রুত যোগাযোগ ও তথ্যের সুযোগের সঙ্গে অন্যের ক্ষতি না করার দায়িত্বও থাকে। কাজ সঠিক কি না শুধু ধরা পড়া দিয়ে নির্ধারিত হয় না। ব্যক্তিগত তথ্য নেওয়ার আগে প্রয়োজন ও সম্মতি, ব্যবহারে সীমিত প্রবেশাধিকার, প্রকাশে সত্যতা ও অধিকার এবং পরে নিরাপদ সংরক্ষণ বা অপসারণ বিবেচনা করতে হয়।
\term{Copyright} সৃজনশীল কাজের অধিকার। \term{Patent} উদ্ভাবনের সঙ্গে সম্পর্কিত অধিকার। অন্যের লেখা, ছবি, সফটওয়্যার বা গবেষণার ফল ব্যবহারে অনুমতি, লাইসেন্স ও প্রযোজ্য শর্ত দেখতে হয়। প্রকাশ্যে পড়া যায় মানেই কপি বা নিজের নামে প্রকাশের অধিকার নয়।
\definition{প্লেজিয়ারিজম}{অন্যের লেখা, ধারণা বা ফল নিজের বলে উপস্থাপন করাকে Plagiarism বলে।}
অধিকার বা লাইসেন্স লঙ্ঘন করে ব্যবহার Copyright Infringement হতে পারে। উৎস লিখেও লাইসেন্স ভাঙা সম্ভব; অনুমতি নিয়ে কপি করেও নিজের বলে দাবি করলে প্লেজিয়ারিজম হয়। নিজের ভাষায় ব্যাখ্যা, যথাযথ স্বীকৃতি ও ব্যবহারের অধিকার---তিনটি বিবেচনা দরকার। অন্যের পুরো কাজ হুবহু নিজের বলে প্রকাশ করা Clone Plagiarism-এর উদাহরণ।
\term{Software Piracy} হলো লাইসেন্সের শর্ত ভেঙে সফটওয়্যার কপি, ব্যবহার বা বিতরণ। বৈধ লাইসেন্স বা উপযুক্ত মুক্ত বিকল্প ব্যবহার করতে হয়। অন্যের তথ্য মুছে বা বিকৃত করাও ক্ষতিকর; তথ্যের অখণ্ডতা রক্ষা জরুরি।
\subsection{অনলাইন অপরাধ ও প্রতারণা}
\begin{compactpoints}
\item \term{Phishing}: বিশ্বস্ত ব্যক্তি বা প্রতিষ্ঠান সেজে ভুয়া বার্তা, লিংক বা ওয়েবসাইটে গোপন তথ্য আদায়ের চেষ্টা। মানুষকে ঠকানোর কৌশল Social Engineering।
\item \term{Spoofing}: পরিচয়, প্রেরকের ঠিকানা বা অন্য শনাক্তকারী তথ্য নকল করে বিশ্বাসযোগ্য সেজে থাকার কৌশল। ফিশিংয়ে এমন পরিচয় নকল করা হতে পারে; দুটি শব্দ একই অর্থ বোঝায় না।
\item \term{Spam}: অনাকাঙ্ক্ষিত ব্যাপক বার্তা। ছাঁকনি, ব্লক ও রিপোর্টের ব্যবস্থা ব্যবহার করা যায়।
\item \term{Cyberbullying}: অনলাইনে অপমান, হুমকি বা হয়রানি। প্রমাণ রেখে রিপোর্ট ও বিশ্বস্ত ব্যক্তির সহায়তা নেওয়া দরকার।
\item \term{Unauthorized Access}: অনুমতি ছাড়া অ্যাকাউন্ট বা ব্যবস্থায় প্রবেশ। নিরাপত্তা পরীক্ষা করতেও অনুমতি ও তার সীমা মানতে হয়।
\end{compactpoints}
Hacker দক্ষ প্রোগ্রামার বা নিরাপত্তা গবেষকও হতে পারেন। White Hat অনুমতি ও নির্ধারিত সীমায় পরীক্ষা করেন; Black Hat ক্ষতিকর বা অননুমোদিত উদ্দেশ্যে কাজ করেন। Grey Hat-এর প্রচলিত বর্ণনায় অনুমতি ছাড়া ত্রুটি খোঁজা থাকতে পারে; ভালো উদ্দেশ্য দাবি অনুমতির বিকল্প নয়। Cracker সাধারণত ক্ষতিকর বা অননুমোদিত নিরাপত্তা ভাঙার প্রসঙ্গে ব্যবহৃত হয়।
\subsection{ক্ষতিকর সফটওয়্যার ও সুরক্ষা}
ক্ষতিকর সফটওয়্যারের সাধারণ নাম Malware। Virus সাধারণত হোস্ট ফাইলে যুক্ত হয়ে ছড়ায়। Worm নিজে ছড়িয়ে অন্য ব্যবস্থায় পৌঁছাতে পারে। Trojan উপকারী সফটওয়্যারের ছদ্মবেশে ক্ষতি করে। Spyware গোপনে তথ্য সংগ্রহ এবং Ransomware তথ্য বা ব্যবস্থা আটকে মুক্তিপণ দাবি করতে পারে।
শক্তিশালী পাসওয়ার্ড, two-factor authentication, হালনাগাদ সফটওয়্যার, সীমিত অনুমতি ও পৃথক ব্যাকআপ ঝুঁকি কমায়। অপরিচিত লিংক বা সংযুক্তি এড়ানো এবং ব্যক্তিগত ছবি, NID, ব্যাংকের তথ্য ও OTP গোপন রাখা দরকার। প্রকাশের আগে তথ্য যাচাই করতে হয়।
যেমন, কলেজের বৃত্তির নামে বার্তায় OTP চাইলে বিশ্বস্ত মাধ্যমে অফিসে যাচাই করতে হবে; ওই লিংকে লগইন করা উচিত নয়। তথ্য দিয়ে ফেললে দ্রুত পাসওয়ার্ড বদলানো, সক্রিয় সেশন প্রত্যাহার ও সেবাদাতাকে জানানো দরকার। সাধারণ অনাকাঙ্ক্ষিত বিজ্ঞাপন এবং গোপন তথ্য নেওয়ার প্রতারণা পৃথক ঘটনা।
\subsection{কম্পিউটার ব্যবহারের দশটি নৈতিক নীতি}
\begin{multicols}{2}\begin{compactenum}
\item কম্পিউটার দিয়ে অন্যের ক্ষতি না করা।
\item অন্যের কাজ ব্যাহত না করা।
\item অনুমতি ছাড়া অন্যের ফাইল না দেখা।
\item কম্পিউটার দিয়ে চুরি না করা।
\item মিথ্যা বা বিভ্রান্তিকর তথ্য না দেওয়া।
\item অধিকার ছাড়া সফটওয়্যার কপি বা ব্যবহার না করা।
\item অনুমতি বা ন্যায্য প্রতিদান ছাড়া অন্যের সম্পদ ব্যবহার না করা।
\item অন্যের মেধাসম্পদ নিজের বলে দাবি না করা।
\item ব্যবস্থা তৈরির সামাজিক ফল বিবেচনা করা।
\item মানুষের মর্যাদা ও অধিকার সম্মান করা।
\end{compactenum}\end{multicols}
\chOnePair{\chOneNative{Ethics}}{তথ্য নেওয়া, ব্যবহার, প্রকাশ ও সংরক্ষণের প্রতিটি ধাপে দায়িত্বশীল সিদ্ধান্ত দরকার।}
\begin{practicebox}
\begin{enumerate}
\item ICT নৈতিকতা কী?
\item Software piracy কেন অনৈতিক?
\item নিরাপদ internet ব্যবহারের চারটি নিয়ম লেখ।
\end{enumerate}
\end{practicebox}
